% META: {"id":"1SPE-PRODSCAL-EX-032","chapitre":"1SPE-PRODUIT-SCALAIRE","type_objet":"exercice","capacites":["1SPE-PRODUIT-SCALAIRE-C4"],"capacites_codes":["C4"],"methodes":["M4"],"parcours":1,"competences":["calculer"],"duree_min":5,"mode_creation":"ex_nihilo","sources_inspiration":[],"fichier_tex":"chapitres/1SPE-PRODUIT-SCALAIRE/exercices/1SPE-PRODSCAL-EX-032.tex","corrige_tex":"chapitres/1SPE-PRODUIT-SCALAIRE/corriges/1SPE-PRODSCAL-CO-032.tex","status":"approved","origin":"generated","status_history":["generated","audited","accepted","approved"],"release_acceptance":"RELEASE_OWNER_BATCH_ACCEPTANCE","acceptance_closure_digest":"sha256:9b3ccf9a81c5520fbb7e03b7b2d3e7bf2057a02834b6d908bf3be5f49f3b553f","acceptance_receipt_digest":"sha256:4fad86f0282e0fa4e0419cca1ad6379ae7af5a280c04a4a90880f90a1c044242"}
% BEGIN-VERIFY
% from sympy import *
% # A(0,0), B(6,0), C(2,4): AH perp BC
% # BC = (-4,4), H on line BC: H = B + t*BC = (6-4t, 4t)
% # AH.BC = 0: (6-4t)*(-4) + 4t*4 = 0 => -24+16t+16t=0 => 32t=24 => t=3/4
% t = Rational(3, 4)
% Hx = 6 - 4*t
% Hy = 4*t
% assert Hx == 3
% assert Hy == 3
% END-VERIFY
\begin{exercice}{1SPE-PRODSCAL-EX-032}{1}{8}
Soit $A(0; 0)$, $B(6; 0)$, $C(2; 4)$. Déterminer les coordonnées du pied $H$ de la hauteur issue de $A$ dans le triangle $ABC$.
\end{exercice}
